%%%%%%%% Project proposal on the MLSys 2024 template %%%%%%%%%%%%%%%%%

\documentclass{article}

\usepackage{microtype}
\usepackage{graphicx}
\usepackage{subfigure}
\usepackage{booktabs}
\usepackage{hyperref}

\newcommand{\theHalgorithm}{\arabic{algorithm}}

% Blind version; switch to [accepted] to show the team names below.
\usepackage[accepted]{mlsys2024}% \usepackage[accepted]{mlsys2024}

\mlsystitlerunning{Globally Coordinated Speculative Decoding}

\begin{document}

\twocolumn[
\mlsystitle{Coordinated Adaptive Speculative Decoding
for Distributed LLM Serving}

\mlsyssetsymbol{equal}{*}

\begin{mlsysauthorlist}
\mlsysauthor{Alvin Wang}{cmu}
\mlsysauthor{Arjun Shah}{cmu}
\end{mlsysauthorlist}

\mlsysaffiliation{cmu}{Carnegie Mellon University, Pittsburgh, PA, USA}

\mlsyscorrespondingauthor{Alvin Wang}{alvinw2@andrew.cmu.edu}

\mlsyscorrespondingauthor{Arjun Shah}{aashah2@andrew.cmu.edu}

\mlsyskeywords{LLM Serving, Speculative Decoding, Expert Parallelism, Scheduling}

\vskip 0.3in

\begin{abstract}
Adaptive speculative decoding is now standard inside a single serving engine. Large deployments, however, run many workers that are coupled, such as data-parallel attention ranks that advance in lockstep with expert parallelism, or verifiers that share a pool of drafters. In these
settings, each worker optimizing its own speculation budget is not globally
optimal, because the slowest rank sets the step time for everyone. We propose
CoSpec, a coordinator that chooses speculation budgets for each request across all
workers jointly, and co-designs request routing with those budgets. We will build
CoSpec on an open-source engine and evaluate it on production-like traces against
per-engine adaptive baselines.
\end{abstract}
]

\printAffiliationsAndNotice{}

%-------------------------------------------------------------------------------
\section{Problem and System Challenges}
\label{sec:problem}

Speculative decoding trades extra compute for fewer sequential steps. At low
load, verifying extra draft tokens is nearly free. At high load, every rejected
token occupies the batch capacity that another request could have used, so a fixed draft
length can add latency compared to plain decoding
\cite{liu2024smartspec,cheng2026dspark}. The best budget therefore depends on both
the request (its acceptance probability) and the system (its load).

Existing adaptive schedulers answer this question for one engine. Modern
deployments are distributed, which raises three challenges these schedulers do not address:

\begin{itemize}
\item \textbf{Lockstep coupling:} In data parallel (DP) attention and expert parallel (EP) Mixture-of-Experts
serving, all DP ranks must finish each step together; idle ranks even run dummy
passes \cite{anyscale2026ep}. A rank's speculation budget changes its batch size
and step time, and the slowest rank gates the group \cite{nanocp2026}. Local
greedy choices waste the slack on lightly loaded ranks and overload heavy ones.
\item \textbf{Routing ignores speculation.} A request's best speculation
budget depends on which rank it runs on. For example, a lightly loaded rank can verify many
draft tokens cheaply, a busy one cannot. In turn, speculation changes how busy
a rank really is, since each request may add several verification tokens per
step. Routers today balance request counts or queue lengths and ignore this,
and speculation schedulers take placement as given, so a high-acceptance request can land on a busy rank where it gets little budget.
\item \textbf{Shared draft capacity.} Disaggregated drafters serving many
verifiers \cite{vllm2026disagg,he2026nebulasd} need a policy for dividing a
finite budget among verifiers with different acceptance and load.
\end{itemize}

Doing this efficiently is a systems problem. The coordinator must decide
within a single decode step (a few ms), use little communication, and never break the
lossless guarantee of speculative decoding, which forbids budget decisions that peek at future draft
tokens \cite{cheng2026dspark}.

%-------------------------------------------------------------------------------
\section{State of the Art}
\label{sec:sota}

Table~\ref{tab:sota} summarizes the closest work
\cite{liu2024smartspec,li2025adaserve,huang2025adaspec,wu2025tetris,chen2026faser,tran2026goodspeed}.
Per-engine methods adapt to
acceptance and load but assume each engine is independent. DSpark explicitly
assumes a load balancer keeps DP ranks balanced so that throughput depends only on
local batch size \cite{cheng2026dspark}. Distributed systems so far either use
heuristic depth rules evaluated on a single 4-GPU node \cite{kumar2026streamserve},
or focus on pooling mechanics rather than budget policy \cite{he2026nebulasd}, or
target edge settings with one verifier \cite{tran2026goodspeed}.

\begin{table}[t]
\caption{Closest prior work. ``Scope'' is where budgets are decided.
Code status is as of Sept.\ 2026; ``--'' means we did not find a public release.}
\label{tab:sota}
\vskip 0.1in
\begin{center}
\begin{footnotesize}
\setlength{\tabcolsep}{4pt}
\begin{tabular}{lll}
\toprule
System & Scope & Code \\
\midrule
SmartSpec & 1 engine & vLLM-based \\
AdaServe & 1 engine, per SLO & -- \\
AdaSpec & 1 engine & public \\
TETRIS & 1 batch & -- \\
FASER & 1 engine & -- \\
DSpark & 1 engine & drafter, ckpts \\
StreamServe & per worker (heuristic) & -- \\
NebulaSD & draft/target pools & -- \\
GoodSpeed & 1 verifier, $N$ drafters & -- \\
\bottomrule
\end{tabular}
\end{footnotesize}
\end{center}
\vskip -0.15in
\end{table}

Open building blocks we can reuse: vLLM and SGLang (SD with EAGLE-3 and MTP,
DP+EP deployment), public EAGLE-3 drafters, and DeepSpec checkpoints and training
code \cite{cheng2026dspark}.

%-------------------------------------------------------------------------------
\section{Proposed Directions}
\label{sec:directions}

Each direction listed below addresses one of the challenges and gaps in current SOTA schedulers. The first direction is to determine the global budget allocation in lockstep ranks. Each rank $r$ has candidate draft tokens with survival probabilities $a_{r,j}$
and a profiled step-time model $T_r(B_r, \mathrm{KV}_r)$. Because the group step
time is $\max_r T_r$, we maximize the total expected accepted tokens divided by
$\max_r T_r$. A greedy ``water-filling'' allocator
admits the highest-value tokens globally until any rank's step time would rise,
so lightly loaded ranks spend their slack on speculation. We will keep decisions
non-anticipating to stay lossless, and piggyback the needed summaries
(a few scalars per rank) on the existing DP synchronization.

The second direction is to implement speculation-aware routing. Extend the decode router to score each
rank by expected goodput after budget allocation, not just queue length,
steering high-acceptance requests (e.g., code) toward ranks with slack.

The third direction is to create a shared draft pool. We can allocate draft/verify capacity
across $N$ verifiers sharing $M$ drafters with the same marginal-value rule.

%-------------------------------------------------------------------------------
\section{Evaluation Plan}
\label{sec:eval}

Our compute is a single NVIDIA V100, so we cannot run real lockstep DP+EP
groups. We therefore use a trace-driven methodology. The GPU supplies
real costs and real acceptance behavior, and a simulator replays them across
many virtual ranks.

\textbf{Measurement}: We run Qwen3-4B in fp16 with released
EAGLE-3 and DSpark drafters \cite{cheng2026dspark}, using vLLM if it supports
the V100 and a minimal PyTorch decode loop otherwise. We (i) profile verification
step time $T(B,\mathrm{KV})$ across batch sizes and context lengths, and
(ii) record, for every drafted token, the drafter's predicted confidence and
whether the target accepted it.

\textbf{Simulation:} A discrete-event simulator models $R\in\{4,8,16,64\}$
lockstep ranks whose group step time is $\max_r T_r$, using the profiled curves
and replaying the recorded acceptance traces. It implements routing and a
shared draft pool.

\textbf{Single-GPU validation:} We execute several virtual ranks' verification
batches back to back on the V100, charge each group step their maximum time,
and check that measured throughput matches the simulator's prediction. This
also measures the allocator's per-step overhead.

\textbf{Workloads:} A code-heavy mix (HumanEval, MBPP) and a chat-heavy mix
(ShareGPT, MT-Bench) to vary acceptance; arrival times from the Azure LLM
inference traces, with deliberately skewed rank loads.

\textbf{Baselines:} No SD; fixed draft length $k=4$; per-rank adaptive
(DSpark-style greedy on each rank); CoSpec with Direction 1 only, and with Direction 1 + Direction 2.

\textbf{Metrics:} Goodput under TPOT SLOs, P50/P99 TPOT, throughput.

%-------------------------------------------------------------------------------
\section{Team and Division of Work}
\label{sec:team}

\begin{itemize}
\item \textbf{Alvin Wang (GPU side):} setting up the decode loop on the V100,
profiling step-time curves, collecting acceptance traces, running the
single-GPU validation, and preparing workloads.
\item \textbf{Arjun Shah (algorithm and simulator side):} building the lockstep
simulator, implementing the allocator (D1), routing (D2) and shared draft
pool (D3), and running the simulated baselines.
\end{itemize}
Both members share benchmarking, analysis, writing, and meet weekly to
check the simulator against GPU measurements.


%-------------------------------------------------------------------------------
\section{Timeline}
\label{sec:timeline}

\begin{table}[h]
\caption{Milestones (Oct.\ 1 -- Dec.\ 12, 2026).}
\label{tab:timeline}
\vskip 0.1in
\begin{center}
\begin{footnotesize}
\setlength{\tabcolsep}{4pt}
\begin{tabular}{ll}
\toprule
Weeks & Milestone \\
\midrule
1--2 & Pick engine; DP+EP + SD running on small MoE \\
3--4 & Profiling on full node; build simulator \\
5--6 & D1 in simulator, then in engine (small MoE) \\
7 & \emph{Midterm:} D1 vs.\ baselines, first full-node runs \\
8--9 & D2 routing; D3 in simulator \\
10 & Final full-node runs; simulated scaling \\
11 & Final report and presentation \\
\bottomrule
\end{tabular}
\end{footnotesize}
\end{center}
\vskip -0.15in
\end{table}

%-------------------------------------------------------------------------------
\section{Infrastructure and Compute Estimate}
\label{sec:infra}

\textbf{GPU:} A single NVIDIA V100 (32\,GB preferred; 16\,GB is workable with
smaller batches). We use it only for measurement: profiling step-time curves,
collecting acceptance traces, and single-GPU validation. The target is
Qwen3-4B in fp16 ($\sim$8\,GB of weights) with released EAGLE-3 and DSpark
drafters, leaving the rest of memory for the KV cache.

\textbf{CPU:} A 16+ core machine with 64\,GB RAM for the lockstep simulator,
trace processing and analysis. The simulator needs no GPU.

\textbf{Storage:} About 200\,GB: model and drafter weights ($\sim$15\,GB),
acceptance traces and step-time profiles ($\sim$50\,GB), and logs and results.

\textbf{GPU time (about 100 V100-hours):}
\begin{itemize}
\item Setup and debugging of the decode loop: 20 h.
\item Profiling $T(B,\mathrm{KV})$ over batch sizes and context lengths: 15 h.
\item Collecting acceptance traces (2 workload mixes $\times$ 2 drafters): 35 h.
\item Single-GPU validation of the allocator against the simulator: 15 h.
\item Reruns and buffer: 15 h.
\end{itemize}

\textbf{CPU time:} Simulated experiments (5 systems $\times$ 2 workloads
$\times$ 4 rank counts $\times$ 3 load levels $\times$ 3 seeds $=$ 360 runs of a
few minutes each) take roughly 50--100 core-hours. We use existing drafter
checkpoints, so no training is required.

\bibliography{example_paper}
\bibliographystyle{mlsys2024}

\end{document}
